\section{The internal advantage depends on the model}
\label{sec:result}


\begin{table*}[t]
\centering
\small
\caption{Every run under the corrected labels and evaluator. $n_+$ and $n_-$ are the labelled failures and correct calls in the scored population, \emph{attention} the better of the per-head profile and LapEigvals, and the last column the internal advantage with its 95\% tool-resampled interval. $^{*}$ re-extracted under the corrected extractor. Other rows relabelled on their original extraction. $^\dagger$ fewer than thirty items of one class.}
\label{tab:main}
\begin{tabular}{llrrccccc}
\toprule
Model & Data & $n_+$ & $n_-$ & probe & attention & confidence & floor & probe $-$ confidence \\
\midrule
Llama-3.2-1B & Glaive & 223 & 357 & 0.957 & 0.947 & 0.579 & 0.540 & $+0.379$ [+0.15, +0.59] \\
Llama-3.2-1B & BFCL & 193 & 307 & 0.861 & 0.815 & 0.661 & 0.652 & $+0.200$ [+0.12, +0.27] \\
Llama-3.2-1B & BFCL-live & 209 & 156 & 0.780 & 0.795 & 0.661 & 0.707 & $+0.119$ [-0.00, +0.24] \\
Llama-3.2-3B & Glaive & 97 & 548 & 0.936 & 0.885 & 0.773 & 0.529 & $+0.163$ [+0.05, +0.29] \\
Llama-3.2-3B & BFCL & 157 & 539 & 0.884 & 0.833 & 0.712 & 0.621 & $+0.172$ [+0.11, +0.23] \\
Gemma-3-1B & Glaive & 158 & 460 & 0.883 & 0.837 & 0.650 & 0.704 & $+0.233$ [+0.11, +0.34] \\
Gemma-3-1B & BFCL & 379 & 135 & 0.949 & 0.903 & 0.707 & 0.866 & $+0.242$ [+0.19, +0.30] \\
\midrule
Qwen3-1.7B & BFCL & 64 & 627 & 0.791 & 0.745 & 0.741 & 0.616 & $+0.051$ [-0.06, +0.15] \\
Qwen3-1.7B$^\dagger$ & Glaive & 10 & 699 & 0.676 & 0.781 & 0.867 & 0.677 & $-0.191$ [-0.65, +0.21] \\
MiniCPM5-2B & BFCL & 67 & 627 & 0.808 & 0.758 & 0.815 & 0.564 & $-0.008$ [-0.13, +0.09] \\
MiniCPM5-2B & BFCL-live & 123 & 472 & 0.747 & 0.688 & 0.798 & 0.543 & $-0.051$ [-0.16, +0.02] \\
MiniCPM5-2B$^\dagger$ & Glaive & 9 & 716 & 0.710 & 0.713 & 0.913 & 0.685 & $-0.203$ [-0.51, +0.01] \\
Qwen3.5-0.8B & BFCL & 91 & 393 & 0.691 & 0.649 & 0.789 & 0.592 & $-0.099$ [-0.18, -0.02] \\
\bottomrule
\end{tabular}

\end{table*}

\cref{tab:main} and \cref{fig:gap} hold the paper's measurement. The \nRunsPowered{} powered runs fall into two groups that do not overlap. On the seven runs of Llama-3.2 and Gemma-3 the internal advantage runs from \gapInternalsMin{} to \gapInternalsMax{}, and its interval excludes zero on \nInternalsCiExcludesZero{} of them, the exception being Llama-3.2-1B on the live categories, where it touches zero. On the four powered runs of Qwen3, MiniCPM5 and Qwen3.5 it runs from \gapConfidenceMin{} to \gapConfidenceMax{}, its interval is above zero on none and below zero on \nConfidenceCiBelowZero{}. The same ordering holds for a probe that needs no call positions at all, the token-level probe pooled over the span: its advantage is \tokGapMinInternals{} to \tokGapMaxInternals{} on the first group and \tokGapMinConfidence{} to \tokGapMaxConfidence{} on the second. In practice, an operator of a Llama-3.2 or Gemma-3 agent who reads only the model's confidence passes most of the wrong calls that a probe would have stopped, and an operator of a Qwen3, Qwen3.5 or MiniCPM5 agent who builds a probe gains at most a few points over the confidence the model already reports.

Runs that share a checkpoint are not independent, so the unit for any test is the checkpoint, of which there are \nCkptInternals{} on each side. Their mean advantages separate completely, Llama-3.2-1B at \ckptGapLlamaOneB{}, Llama-3.2-3B at \ckptGapLlamaThreeB{}, Gemma-3-1B at \ckptGapGemma{} against Qwen3-1.7B at \ckptGapQwenThree{}, MiniCPM5-2B at \ckptGapMiniCpm{} and Qwen3.5-0.8B at \ckptGapQwenThreeFive{}. Three against three is the smallest attainable two-sided rank-test value, $p=\ckptExactFloor{}$, so the separation is reported as a description of these six checkpoints and not as a test, and \cref{sec:discussion} says what would test it.

Three properties of the model set do not separate the sides on their own. Size: the first side holds checkpoints of one and three billion parameters and the second of 0.8, 1.7 and two billion, and within Llama-3.2 the probe improves from 1B to 3B, as \citet{yeats2026calls} report across sizes, without the advantage changing sign. Call format: Qwen3 writes JSON like Llama-3.2 and Gemma-3 and sits on the other side. Attention design: MiniCPM5 is a dense stack like Llama-3.2, and Gemma-3 attends locally in all but four layers yet behaves like Llama-3.2. Release date comes closest, with every first-side checkpoint released by March 2025 and every second-side one from April 2025 on, and \cref{sec:discussion} returns to what that could mean.

\paragraph{Replication.} Llama-3.2-1B on BFCL was extracted twice, once before and once after the extraction repairs of \cref{sec:setup}, so the two runs differ in their generations as well as in their features. The internal advantage reads \replGapStale{} on the earlier extraction and \replGapClean{} on the later one, a drift of \replGapDelta{} that is inside the interval of either. Differences between detectors smaller than that drift are reported in the tables but not interpreted.
